% ============================================================================
%  qicert — Phase-1 Concept Proposal (portal submission document)
%  Written to Phase-1-Submission-Guidelines-VF.md Sec. 4.3: max 6 pages, 10pt
%  min, seven required elements in order. Companion artifacts: the technical
%  report (docs/technical-report-v2.pdf) and the public repo. Every number is
%  measured and traced in the report's Appendix A; nothing here is projected
%  or aspirational unless explicitly labeled as a Phase-2 target.
% ============================================================================
\documentclass[10pt]{article}
\usepackage[a4paper,margin=2.2cm]{geometry}
\usepackage{amsmath,amssymb}
\usepackage{booktabs}
\usepackage{graphicx}
\usepackage{url}
\urlstyle{same}
\usepackage[hidelinks]{hyperref}
\usepackage{enumitem}
\setlist{nosep,leftmargin=1.4em}
\graphicspath{{figures2/}}
\setlength{\parskip}{0.35em}
\setlength{\parindent}{0pt}

\newcommand{\norm}[1]{\left\lVert #1 \right\rVert}

\begin{document}

\begin{center}
{\LARGE\bfseries qicert: Certified Compression of Vision--Language--Action Models}\\[2pt]
{\large with Exact Lipschitz Bounds and a Refusing Runtime Guard}\\[6pt]
Phase-1 Concept Proposal --- 2026 Global Quantum + AI Challenge\\
Problem statement: \emph{Quantum-Enhanced Vision-Language-Action Models in
Autonomous Driving and Robotics Applications} (Volkswagen Group)\\
Code: \url{https://github.com/ZakLr/qicert} \quad$\cdot$\quad
Technical report: \texttt{docs/technical-report-v2.pdf} (7~pages, in repo)
\end{center}

% ---------------------------------------------------------------------------
\section*{1. Problem Framing}
Vision--language--action models (VLAMs) for driving and robotics are 7B--10B
parameter systems that must eventually run on embedded controllers under
functional-safety regimes (ISO~26262, IEC~62061). The challenge statement
names the compound constraint precisely: shrink the model \emph{and} retain
verifiable behavior on long-tail inputs. Classical compression solves the
first half only. Quantization, pruning, and distillation reduce footprint, but
none of them yields a statement of the form ``under any input perturbation of
size $\delta$, the action cannot leave this set'' --- which is the artifact a
safety engineer would need before certifying deployment on long-tail sensor
distributions.

Quantum or quantum-inspired advantage enters where structure, not scale, does
the work. Our QI component is tensor-network structure (tensor-train
decomposition, the MPS/TTNS family the statement highlights). Its decisive
property is not size reduction --- at the ratios we tested it loses the
accuracy contest to INT8 and we report that plainly --- but that the
compressed weights \emph{carry an exact per-layer Lipschitz constant}
computable from the compressed cores alone, with no reference to data.
Quantized weights admit no comparable exact bound, because quantization error
is data-dependent. Structure is therefore the classical-side object that makes
an exact deployment certificate computable, chainable, and enforceable at
runtime --- and it is structurally unavailable to the accepted baseline. This
is the specific, narrow sense in which the quantum-inspired method provides
the credible advantage the guidelines ask for; we deliberately claim nothing
broader.

% ---------------------------------------------------------------------------
\section*{2. Technical Approach}
\textbf{Paradigm: quantum-inspired (tensor-network), classical pipeline,
no quantum hardware required or assumed.} The statement's simulation clause
covers us explicitly; we state resources honestly in Section~3.

The pipeline wraps an existing VLA at four points. A fine-tuning stage
establishes the reference policy and freezes a hash-verified evaluation split.
The compression stage decomposes each linear layer of the language backbone
into a tensor-train network (TT-SVD); a closed-form low-rank residual folds
the truncation loss back into the operator so that accuracy survives
compression as well as the budget allows, with the certificate cost known
exactly (triangle inequality on the operator norm). The certificate stage
chains the exact per-layer bounds into a model-level perturbation ball and
signs the artifact into a hash manifest. The deployment stage is a runtime
guard that verifies the manifest at load time and, at every inference step,
refuses any action outside the certified ball.

Three measured components complete the stack. First, a commuting-Pauli
interaction compiler evaluates cross-modal interaction operators exactly over
pruned families; its pruning predictions track measurement at median
$r = 0.975$ over twelve independent operator tables, with compiler identity
verified to $2.3\times10^{-14}$ relative error. Second, a rare-event
estimation layer races three estimators on a known-ground-truth problem
($p^{*}=1.37\times10^{-4}$): an iterated quantum amplitude estimation schedule
(simulated; the estimator sees only measurement counts, as it would on
hardware) reaches naive-Monte-Carlo interval width with
$\sim$1{,}587$\times$ fewer queries, and all arms cover the truth at budget.
Third, a conformal layer calibrates violation thresholds whose empirical
coverage lands inside the finite-sample band at all three tested levels
($\alpha = 0.01/0.05/0.10$).

\textbf{Proof status.} The layer bound $\norm{W}_{2} \le \prod_k
\norm{G^{(k)}}_{3}$ is exact algebra on the TT factorization; the certificate
chain and the guard's soundness/completeness properties are machine-checked
by SMT falsification (Z3 finds no counterexample to accepting an out-of-ball
action). A guard demo transcript shows: clean load SERVE, out-of-ball action
REFUSE, tampered-certificate REFUSE, NaN-bound REFUSE, missing-manifest REFUSE.

\textbf{Working prototype.} This is not a design on paper. The repository
contains 80 passing tests, a benchmark harness, and measured results on a
fine-tuned 0.5B VLA (LIBERO-spatial, 114{,}497 scored action tokens,
hash-frozen split): reference accuracy 0.4468; calibrated INT8 at
$1.997\times$ compression, 0.4466 (the honest classical comparator);
tensor-compressed routes at $2.5$--$4.1\times$ with sound certificates on
every one of 168 compressed layers across every configuration tested. The
accuracy of the repaired tensor route under the pre-registered bar is the
open question Section~5 turns into the Phase-2 validation plan; the negative
result at the tested configurations is reported in full in the technical
report rather than hidden.

% ---------------------------------------------------------------------------
\section*{3. Feasibility and Resource Requirements}
Feasibility is demonstrated, not projected: every component above has already
run end to end on consumer hardware. The full experimental program to date
consumed roughly 15~GPU-hours on a single laptop-class GPU (NVIDIA RTX~5060,
8~GB), including fine-tuning (1{,}200 steps), all compression arms, full-split
evaluations, the estimator fixtures, and latency profiling. Peak training
memory was 4.8~GiB. No proprietary data, cloud credits, or quantum hardware
are involved at any stage; LIBERO is the statement's own recommended robotics
dataset; simulation environments and versions are pinned in the repository.
Phase-2 scale-up (below) is sized for the same single-GPU discipline, with
multi-seed replication the main added cost. The honest constraint we declare:
tensor-network contraction on a CPU edge profile costs $3$--$15\times$ the
dense fp16 kernel time per layer at these shapes --- within the 100~ms budget
on measured paths, but a real margin cost that Phase~2 must optimize rather
than wish away.

% ---------------------------------------------------------------------------
\section*{4. Expected Impact}
For the enterprise use case, the deliverable changes what deployment approval
means. A fleet operator today ships a compressed model on the strength of
benchmark scores; with qicert-style artifacts they would ship a signed
certificate, a guard that refuses out-of-ball actions, and an audited
statistical tail --- the difference between ``we tested it'' and ``it is
bounded, and the machine enforces the bound.'' The design tool built on our
degradation measurements points the same direction: predict the largest safe
compression \emph{before} committing to full evaluation, turning post-hoc
luck into engineering process.

Quantitative targets, anchored to what already runs: the reference VLA at
0.4468 accuracy and the INT8 comparator at $1.997\times$ define the bar the
tensor route must clear; certificates already hold at $4.1\times$ (sound
$168/168$); the estimator layer already delivers the $\sim$1{,}587$\times$
query reduction. A successful Phase-2 demonstrates a certified configuration
that is simultaneously (a) at or beyond $2\times$ compression, (b) within the
pre-registered accuracy band (reference $-0.05$), and (c) certificate-sound
with guard enforcement --- or, failing that, the certified sub-$2\times$
regime with a validated predictor that says exactly where the accuracy wall
sits, which is itself deployable information the quantized baseline cannot
produce.

% ---------------------------------------------------------------------------
\section*{5. Validation Plan}
The Phase-2 sprint validates three claims, each with metrics fixed in advance
(pre-registered in the repository before the runs they judge) --- this is the
strongest differentiator we can offer: the gates are committed, mechanical,
and auditable.

\textbf{Capability claim (the live experiment).} Configuration search over
the repaired tensor-compression route --- activation-weighted residual
fitting, mixed quantized/tensor allocation, and sub-$2\times$ targets where
INT8 is byte-capped --- with each promising configuration confirmed on the
full 6{,}496-batch frozen protocol. Metrics: compression ratio, token-level
action accuracy, pre-registered gate (ratio $\ge 2\times$ AND accuracy
$\ge 0.3968$), mean $\pm$ standard deviation over 3 independent seeds
(seeds 1--2 are queued at submission time; seed 0 measured).

\textbf{Safety claim.} Metrics: certificate soundness (currently $168/168$
across all configurations), guard refusal behavior (currently 1.0 on all four
tampering scenarios plus in/out-of-ball cases), estimator coverage of known
ground truth (currently 1.0 at budget for all arms), conformal coverage
within finite-sample bands (currently true at all three levels).

\textbf{Design-tool claim.} The degradation predictor: fit weight-space
features to measured agreement, report leave-one-out $R^{2}$ and
predicted-vs-measured maximum safe ratio. Success: predicted safe ratio
within a stated tolerance of measurement on held-out points.

Integration milestones for the sprint, in order: multi-seed replication;
degradation-curve measurement and predictor validation; mixed-allocation
search to the gate; edge-profile latency optimization for the tensor
contraction (the known margin cost); and a hardware-pathway write-up
estimating qubit count, circuit depth, and shot budget for the amplitude
estimation component, which the challenge lists as a secondary objective and
for which our simulation fixtures already define the interface.

% ---------------------------------------------------------------------------
\section*{6. Hybrid / Cross-Domain Integration}
The design is hybrid by construction, and the division of labor is the point.
Classical deep-learning infrastructure does the training and the heavy
evaluation (PyTorch/CUDA on consumer GPU). The quantum-inspired tensor layer
sits inside the model, where its algebraic structure is load-bearing: it is
what makes the exact bound computable at all. Quantum algorithmic structure
appears as the amplitude-estimation schedule for the statistical tail ---
executed in simulation against the same measurement-count interface it would
have on hardware, so the estimator layer is hardware-swappable without
redesign. Formal methods (SMT falsification via Z3, conformal calibration,
scenario optimization) audit the whole stack mechanically rather than by
argument. Cross-domain transfer is native to the method: the certificate
machinery is backbone-agnostic and dataset-agnostic, so the same pipeline
applies unchanged to the autonomous-driving track (the latency profile we
measure is a CPU edge profile, and the challenge's 100~ms real-time constraint
is the same in both tracks); the robotics-track instantiation here was chosen
because its accepted baseline is reproducible on our hardware.

% ---------------------------------------------------------------------------
\section*{7. Team Capability}
The team combines the three disciplines this proposal straddles.
The lead (Zaki) built and operates the full pipeline end to end: the
fine-tuning harness, the tensor-compression kernel stack with its certificate
algebra, the benchmark/ledger infrastructure with pre-registered gates, and
the formal-methods integration --- evidenced by the public repository's
measured results, 80-test suite, and reproducible artifact chain. Working
practice follows the discipline the reviewers will look for: every reported
number resolves to a run directory with environment snapshot; evaluation
splits are frozen and hashed before experiments; negative results are
reported first-class (the technical report's headline compression finding is
an honest NO-GO under the pre-registered bar, with the failure characterized
rather than framed). AI research agents were used as engineering leverage
under the lead's direction; every claim is traceable to committed artifacts,
and the experiment log records what was run, when, and why for every entry.

% ---------------------------------------------------------------------------
\section*{Enclosures}
\begin{itemize}
\item \textbf{Public repository:} \url{https://github.com/ZakLr/qicert}
--- README quickstart, pinned environment, 80-test suite, benchmark harness,
experiment log, and every artifact behind the numbers above.
\item \textbf{Technical report (7 pages):} \texttt{docs/technical-report-v2.pdf}
in the repository --- method, theory, measured results incl.\ the negative
capability result, ablation isolating the QI component, limitations, and a
claim-to-artifact provenance map (Appendix~A) with seven artifact-generated
figures.
\item \textbf{Team profile:} submitted via the portal form (names, roles,
affiliations, prior quantum/QI and domain experience) per guidelines \S4.1.
\end{itemize}

\end{document}
